% FI-EDL Neurocomputing journal extension — manuscript
% Maintained by fi-edl-writer. Numerical claims must trace to
% results/baseline_comparison_report.md.

\documentclass[review,5p,times]{elsarticle}

\usepackage{amsmath, amssymb, amsthm}
\usepackage{booktabs}
\usepackage{graphicx}
\usepackage{hyperref}
\usepackage{cleveref}
\usepackage{microtype}
\usepackage{tabularx}
\usepackage{xcolor}
\usepackage{pifont}
\usepackage{amssymb}
\usepackage{multirow}
\usepackage[ruled,vlined]{algorithm2e}
\newcommand{\todo}[1]{\textcolor{red}{\textbf{[TODO: #1]}}}

\journal{Neurocomputing}

\begin{document}

\begin{frontmatter}

\title{Spectral-Normalized, KL-Free Evidential Deep Learning:\\
A Simpler Recipe for Calibrated Uncertainty and OOD Detection}

% \author[1]{\todo{first author}}
% \author[1]{\todo{corresponding author}}
% \affiliation[1]{organization={\todo{institution}}}

\begin{abstract}
Evidential deep learning (EDL) places a Dirichlet head on a deterministic
network and reads calibration and out-of-distribution (OOD) signal from the
predicted concentration, an appealingly cheap alternative to ensembles and
Bayesian heads. A recent variant, FI-EDL, additionally proposes a
Fisher-information controller that promises to adapt the KL regularization
strength sample-by-sample. We ask a basic question: what \emph{actually}
drives FI-EDL's strong calibration, and can it be exploited more directly?
We first show that the FIM controller fails to deliver its promised
per-sample adaptivity. The Fisher trace's per-sample standard deviation
is only $6.5\%$ of its mean on CIFAR-10 ($1.6\%$ on MNIST), so the gate
has almost no per-sample signal to adapt to. The default constants
$\beta\!=\!\gamma\!=\!1$, applied to a trace of $\mathcal{O}(K)$, then
saturate the gate to a near-constant scalar that is dataset-dependent --
on MNIST $\overline{\lambda}\!\approx\!4\!\times\!10^{-7}$ leaves the KL
term three orders of magnitude below the fit term, while on CIFAR-10
$\overline{\lambda}\!\approx\!4\!\times\!10^{-4}$ produces an effective
KL contribution comparable to the fit ($\sim\!70\%$); neither matches the
adaptive behaviour the controller was designed for. We then identify a deliberately minimal
replacement: a Brier-only fit, backbone spectral normalization for
natural-image scale, and total evidence $\alpha_0$ as the OOD score. This
three-component recipe attains the best calibration on CIFAR-10 across
two backbones: on VGG-16 the recipe's $\text{ECE}\!=\!3.93\!\pm\!0.45\%$
is best-mean (paired-tied with DAEDL, $p\!=\!0.55$), and on ResNet-18
the recipe's $\text{ECE}\!=\!3.15\!\pm\!0.57\%$ is the unique
statistically significant best, dominating each of DAEDL, F-EDL, and
Re-EDL ($p\!\le\!0.04$). On out-of-distribution detection the recipe
achieves best-mean CIFAR-10$\to$SVHN AUROC on both backbones (VGG-16
$92.99\!\pm\!1.77$, ResNet-18 $93.02\!\pm\!1.56$; paired-tied with
Re-EDL on each), uniquely significant CIFAR-10$\to$CIFAR-100 AUROC on
VGG-16 ($86.32\!\pm\!0.20$, $p\!<\!0.01$ over each baseline), and
uniquely significant CIFAR-10$\to$DTD (Describable Textures, a
third OOD set) AUROC of $91.09\!\pm\!0.28$ ($p\!<\!0.05$ over each
baseline), with the one non-best metric being CIFAR-10$\to$CIFAR-100
on ResNet-18, where Re-EDL leads by $0.86$ percentage points.
On the harder CIFAR-100 benchmark (ResNet-18, matching F-EDL's setup)
the recipe is paired-tied with Re-EDL across every metric on four
OOD directions (SVHN, CIFAR-10, TinyImageNet, DTD), while DAEDL's
post-hoc density head fails to train at $K\!=\!100$, an analogous
collapse to the one independently reported in F-EDL's paper.
Three honest ablations support the design: re-enabling the gate with the
original full-$\alpha$ KL destroys calibration (ECE jumps from $5.19\%$ to
$71.62\%$); a target-masked KL controller preserves calibration but adds
nothing; and spectral normalization \emph{harms} calibration on a small
MNIST ConvNet (ECE $1.70\!\to\!13.60\%$), so the recipe recommends
SN-off in small-backbone regimes. The recipe trails DAEDL by a
statistically significant $1.29$ percentage points on misclassification
AUROC, which we attribute to its post-hoc density head and discuss as a
known limitation. The take-away is an Occam's-razor finding for modern
EDL: stripping away the FIM controller, post-hoc density models, and
auxiliary heads of recent baselines yields a simpler recipe that is
competitive or better on the studied benchmarks across calibration and
OOD detection, with one acknowledged loss on misclassification.
\end{abstract}

\begin{keyword}
evidential deep learning \sep Dirichlet networks \sep uncertainty calibration
\sep out-of-distribution detection \sep spectral normalization
\end{keyword}

% --- Elsevier "Highlights" — 5 bullets, each <=85 chars. ----
\begin{highlights}
  \item FI-EDL's Fisher-information gate fails to adapt: trace std is 6.5\% of mean.
  \item Brier-only fit + backbone spectral norm + $\alpha_0$ score: a minimal EDL recipe.
  \item Significant CIFAR-10 ECE win on ResNet-18 ($3.15$, $p{\le}.04$); best-mean on VGG-16.
  \item CIFAR-10$\to$SVHN AUROC 92.99 (VGG) / 93.02 (ResNet), best mean on both backbones.
  \item Honest limitation: spectral norm regresses MNIST ECE from 1.70 to 13.60\%.
\end{highlights}

\end{frontmatter}

% =========================================================================
\section{Introduction}\label{sec:intro}

Modern deep classifiers are accurate but systematically miscalibrated, and
their softmax confidence is a poor signal for whether an input is
out-of-distribution~\citep{guo2017calibration,hendrycks2017baseline}.
A practical uncertainty estimator should deliver three things in a single
deterministic forward pass: a well-calibrated class probability, a usable
score for OOD detection, and a usable score for separating correct from
incorrect predictions. Ensembles and approximate-Bayesian heads do achieve
this, but at a cost in compute and engineering effort that scales with the
backbone.

Evidential deep learning (EDL) is an appealingly cheap alternative. By
re-interpreting the network output as the concentration parameters of a
Dirichlet distribution over class probabilities, EDL exposes a single
``total evidence'' scalar that captures epistemic
uncertainty~\citep{sensoy2018edl,malinin2018prior,amini2020evidential}.
A line of follow-up work has reshaped the training loss around this
parameterization: I-EDL reweights per-class terms with a Fisher-information
matrix~\citep{deng2023iedl}; R-EDL relaxes nonessential design choices and
treats the prior weight as a free hyperparameter~\citep{chen2024redl};
Re-EDL drops the KL-to-uniform regularizer entirely
\citep{chen2025reedl}; DAEDL augments the head with a
post-hoc Gaussian mixture density on top of a spectrally normalized
backbone~\citep{yoon2024daedl}; and Flexible EDL (\citealp{yoon2025fedl},
abbreviated ``F-EDL'' in its venue) replaces the Dirichlet with a Flexible
Dirichlet, equips the backbone and the $\alpha$-head with spectral
normalization, and adds two auxiliary sub-heads predicting the mean $p$
and a temperature $\tau$. Each step adds machinery, but the question of
which pieces are actually doing the work has not been answered head-on.

The Fisher-information EDL paper~--~henceforth \emph{FI-EDL}, abbreviated
this way throughout to avoid the published collision with Flexible EDL~--~is
a particularly clean test case. Its proposed mechanism is an adaptive gate
$\lambda(\alpha)\!=\!\beta\exp\!\bigl(-\gamma\,\mathrm{tr}(I(\alpha))\bigr)$
that multiplies the KL-to-uniform term, so that confident regions of input
space (high Fisher trace) are regularized less and uncertain regions more.
Empirically the method calibrates well, but a quick instrumentation tells
a different story. With the published $\beta\!=\!\gamma\!=\!1$, the Fisher
trace of a Dirichlet on $K\!=\!10$ classes sits at
$\mathrm{tr}(I(\alpha))\!\approx\!14.5\!=\!\mathcal{O}(K)$ throughout
training. The gate $\lambda\!=\!\exp(-\mathrm{tr}\,I(\alpha))$ collapses
the per-sample trace through an exponential, and the empirical
$\overline{\lambda}$ depends on the trace's per-sample distribution: on
MNIST $\overline{\lambda}\!\approx\!4\!\times\!10^{-7}$ and the effective
KL contribution is three orders of magnitude below the fit term, while
on CIFAR-10 $\overline{\lambda}\!\approx\!4\!\times\!10^{-4}$ and the
effective KL contribution is roughly $70\%$ of the fit term
(\autoref{tab:gate}). More striking is the second number: across the
training set, the per-sample standard deviation of
$\mathrm{tr}(I(\alpha))$ is only $6.5\%$ of its mean on CIFAR-10 and
$1.6\%$ on MNIST, so the gate has almost no per-sample variation to
adapt to even after normalisation. The strong calibration of FI-EDL is
real, but it is not produced by an \emph{adaptive} controller -- the
controller acts as a (mis-tuned) constant KL coefficient, fully
saturated on MNIST and moderate on CIFAR-10.

This paper takes that diagnosis as a starting point and asks: what
\emph{actually} drives FI-EDL's calibration advantage, and can it be
exploited more directly? Our answer is a deliberately minimal recipe with
three components: (i) train with the Brier-like fit term alone, dropping
the KL regularizer that the gate had silently nullified anyway;
(ii) apply spectral normalization to the backbone (only) in
natural-image regimes, following the distance-aware Lipschitz argument of
SNGP and DUQ~\citep{liu2020sngp,vanamersfoort2020duq}; and (iii) use the
total Dirichlet evidence $\alpha_0$ as the OOD score, matching the
canonical evidential vacuity uncertainty~\citep{sensoy2018edl} but with a
confidence-like orientation that is a drop-in replacement for the maximum
softmax probability~\citep{hendrycks2017baseline}. Each component
corresponds to \emph{removing} a piece of recent EDL machinery: the FI-EDL
gate, the DAEDL post-hoc GMM head, and the Flexible-EDL auxiliary
$(p,\tau)$ sub-heads on a non-Dirichlet family. The recipe adds no
parameters over plain EDL, requires no auxiliary density model, and trains
in one pass.

Empirically, the recipe is competitive or better than every published EDL
baseline we re-trained under a matched protocol on CIFAR-10 with a VGG-16
backbone. The simple count is three of five metrics first place, one
second, one top-tier; once read through paired $t$-tests and bootstrap
confidence intervals (\autoref{sec:main}, the supplementary material),
the honest statement is: the CIFAR-10$\to$CIFAR-100 OOD AUROC of
$86.32\!\pm\!0.20$ is the unique statistically significant first place
($p\!<\!0.01$ over DAEDL, F-EDL, and Re-EDL); the calibration ECE of
$3.93\!\pm\!0.45\%$ is best-mean but paired-tied with DAEDL ($p\!=\!0.55$);
the CIFAR-10$\to$SVHN AUROC of $92.99\!\pm\!1.77$ is best-mean but paired-
tied with Re-EDL and F-EDL. On misclassification AUROC, DAEDL retains a
significant $1.29$-pp lead. We report all four verdicts as stated rather
than rounding to a slogan.

We also report what does \emph{not} work. Re-enabling the gate with the
original full-$\alpha$ KL collapses calibration on CIFAR-10 from
$5.19\%$ ECE to $71.62\%$ at unchanged accuracy, the unmistakable
signature of severe underconfidence
(\autoref{sec:abl-gate}). A target-masked KL avoids the collapse, but
yields no metric improvement over the KL-free baseline because the
Fisher trace's near-zero per-sample variance leaves the adaptive
controller no signal to act on (\autoref{sec:abl-masked}). And spectral
normalization, which helps every CIFAR-10 metric, \emph{regresses} MNIST
calibration from $1.70\%$ to $13.60\%$ ECE because a small ConvNet under
SN is over-constrained and the $\alpha_0$ score orientation inverts on
FMNIST (\autoref{sec:abl-sn}). The recipe therefore prescribes SN-on for
natural-image backbones and SN-off for small-backbone settings, with the
regression flagged rather than buried.

\paragraph{Contributions.}
\begin{enumerate}
  \item \emph{Diagnosis of the FIM controller's empirical inertia.} We
        instrument FI-EDL during matched re-training on MNIST and CIFAR-10
        and show that the Fisher trace sits at $\mathcal{O}(K)$ with a
        per-sample standard deviation of only $6.5\%$ of its mean, so the
        gate $\lambda\!=\!\exp(-\mathrm{tr}(I(\alpha)))$ saturates at
        $\sim\!5\!\times\!10^{-7}$ and its effective KL contribution is
        five to seven orders of magnitude below the fit
        term (\autoref{sec:diagnosis}, \autoref{tab:gate}).
  \item \emph{A minimal Brier+SN+$\alpha_0$ recipe and its mechanistic
        rationale.} We propose and justify a three-component design that
        removes the FIM gate, the DAEDL density head, and the Flexible-EDL
        $(p,\tau)$ sub-heads, keeping only what the diagnosis implies is
        load-bearing (\autoref{sec:recipe}, \autoref{sec:why}).
  \item \emph{A new state-of-the-art on CIFAR-10 across two backbones
        with honest statistical framing.} On the matched re-train (5
        seeds, 6 EDL baselines) the recipe attains best-mean ECE on
        both VGG-16 ($3.93$, paired-tied with DAEDL, $p\!=\!0.55$) and
        ResNet-18 ($3.15$, \emph{uniquely significant} over each
        baseline, $p\!\le\!0.04$); best-mean CIFAR-10$\to$SVHN AUROC on
        both backbones (paired-tied with Re-EDL); and uniquely
        significant CIFAR-10$\to$CIFAR-100 AUROC on VGG-16
        ($86.32\!\pm\!0.20$, $p\!<\!0.01$ vs.\ each baseline). The one
        non-best metric is CIFAR-10$\to$CIFAR-100 on ResNet-18, where
        Re-EDL leads by $0.86$ percentage points ($p\!<\!0.01$)
        (\autoref{sec:main}, \autoref{tab:main-resnet}). On
        additional CIFAR-10 OOD sets adopted by Re-EDL --- DTD
        \citep{cimpoi2014describing} and GTSRB
        \citep{stallkamp2012gtsrb} --- the recipe is best-mean on both
        (uniquely significant on DTD, paired-tied with Re-EDL on GTSRB)
        (\autoref{tab:main-extended-ood}). On the harder CIFAR-100
        benchmark (ResNet-18, matching F-EDL's Appendix~E.2) it is
        paired-tied with Re-EDL across accuracy, ECE, and four OOD
        directions (SVHN/CIFAR-10/TinyImageNet/DTD), while DAEDL's
        post-hoc density head fails to train under $K\!=\!100$
        (\autoref{tab:main-cifar100}).
  \item \emph{Honest negative results.} We document three: (i) restoring
        the gate with full-$\alpha$ KL destroys calibration
        ($5.19\!\to\!71.62\%$ ECE, \autoref{sec:abl-gate}); (ii) a target-
        masked KL preserves calibration but yields no metric gain because
        the Fisher controller has no signal
        (\autoref{sec:abl-masked}); and (iii) spectral normalization
        regresses MNIST calibration ($1.70\!\to\!13.60\%$ ECE,
        \autoref{sec:abl-sn}), so the recipe recommends SN-off for small
        backbones.
  \item \emph{Cross-baseline score audit.} Switching the OOD score from
        $\max_c\hat p_c$ to total evidence $\alpha_0$ shifts the
        leaderboard meaningfully: the matched-protocol FI-EDL baseline on
        CIFAR-10$\to$SVHN moves from $89.89$ AUROC with $\max_c\hat p_c$
        to $91.72$ with $\alpha_0$ at no retraining cost. The same
        re-evaluation reorders the published leaderboards of several
        other evidential methods (\autoref{sec:score-ablation}),
        supporting $\alpha_0$ as the right default OOD score for
        evidential methods.
\end{enumerate}

\paragraph{Outline.}
\autoref{sec:related} positions FI-EDL+SN next to DAEDL and Flexible EDL
in a head-to-head table that defuses the scoop risk on spectral
normalization. \autoref{sec:background} reviews the Dirichlet/EDL
formulation. \autoref{sec:method} contains the diagnosis
(\autoref{sec:diagnosis}), the recipe (\autoref{sec:recipe}), and the
component-by-component rationale (\autoref{sec:why}). \autoref{sec:exp}
reports main results with paired statistical tests. \autoref{sec:ablations}
presents the three negative results above. \autoref{sec:discussion}
discusses when SN helps and \autoref{sec:conclusion} concludes.

% =========================================================================
\section{Related Work}\label{sec:related}

\subsection{Evidential deep learning}\label{sec:rw-edl}
Evidential deep learning places a Dirichlet (or Normal--Inverse--Gamma for
regression) prior on the predictive distribution and trains a deterministic
network to output its parameters in a single forward pass. Sensoy
\textit{et al.}~\citep{sensoy2018edl} introduce the classification
formulation with a Brier-like fit and a KL-to-uniform regularizer on the
target-masked concentration. Malinin and Gales~\citep{malinin2018prior}
develop the parallel Prior Networks line, which trains by minimizing a KL
between predicted and target Dirichlets built from sharp in-distribution
and flat synthetic OOD samples. Amini \textit{et al.}~\citep{amini2020evidential}
extend the program to regression with a Normal--Inverse--Gamma prior.
Deng \textit{et al.}'s I-EDL~\citep{deng2023iedl} reweights the per-class
EDL loss with the Fisher Information Matrix of the Dirichlet so that
low-evidence classes receive larger gradient. Chen \textit{et al.}'s
R-EDL~\citep{chen2024redl} identifies the fixed prior weight and the
variance-minimizing term as nonessential and treats the prior weight as a
tunable hyperparameter; its journal extension Re-EDL~\citep{chen2025reedl}
goes further and \emph{drops the KL regularizer entirely}, making it the
closest published precedent for the KL-free leg of our recipe.

\subsection{Concurrent work: spectral norm meets evidential
heads}\label{sec:rw-concurrent}
Two recent works -- DAEDL~\citep{yoon2024daedl} and Flexible EDL
\citep{yoon2025fedl} -- combine spectral normalization with Dirichlet (or
Flexible-Dirichlet) heads, and overlap with our recipe enough to require a
head-to-head positioning. The naming is unfortunate: Flexible EDL is
abbreviated ``F-EDL'' in its venue, which collides with an earlier internal
shorthand for Fisher-information EDL; we use \emph{FI-EDL} throughout for
the Fisher-information method and reserve F-EDL for Flexible EDL.

DAEDL~\citep{yoon2024daedl} is, to our knowledge, the first published EDL
method to apply spectral normalization to the feature extractor, and uses
the resulting bi-Lipschitz feature map to fit a class-conditional GMM at
test time whose density re-weights the predicted logits. F-EDL
\citep{yoon2025fedl} replaces the Dirichlet with a Flexible Dirichlet
mixture, applies spectral normalization to both the backbone $f_\theta$
and the $\alpha$-head $g_{\phi_1}$, adds two auxiliary sub-heads predicting
the mixture mean $p$ and a temperature $\tau$, and replaces the KL penalty
with a Brier-score regularizer on $p$. Both share the SN+evidential thesis
with FI-EDL+SN, but each \emph{adds} machinery -- post-hoc density in
DAEDL, a different distribution family plus auxiliary heads in F-EDL.

Our positioning, summarized in \autoref{tab:rw-concurrent}, is an
\emph{Occam's-razor simplification}: we keep the plain Dirichlet, drop
the KL term, drop any post-hoc density, drop the auxiliary heads, and
apply spectral normalization only to the backbone. \autoref{sec:main}
shows that this minimal recipe matches or beats both methods on the
studied CIFAR-10 metrics. A 2026 concurrent line, GEM-FI
\citep{mohammed2026gemfi}, shares the ``Fisher'' label but takes a
mechanistically distant route: it predicts a mixture over evidential
heads with an energy-based gate, neither dropping the KL term nor
applying SN; we therefore treat it as orthogonal to the present design.

\begin{table}[t]
  \centering
  \caption{Side-by-side positioning of FI-EDL+SN against the two
    concurrent SN+evidential methods. The right four columns indicate
    whether each method adds the listed component on top of plain EDL.
    Our recipe is the unique entry that adds nothing.}
  \label{tab:rw-concurrent}
  \small
  \begin{tabular}{lccccc}
    \toprule
    Method & Backbone SN & Drops KL & Extra heads (+ SN) &
      Post-hoc density & Distribution \\
    \midrule
    Plain EDL~\citep{sensoy2018edl} & \ding{55} & \ding{55} & \ding{55} & \ding{55} & Dirichlet \\
    Re-EDL~\citep{chen2025reedl}    & \ding{55} & \checkmark & \ding{55} & \ding{55} & Dirichlet \\
    DAEDL~\citep{yoon2024daedl}     & \checkmark & \ding{55}  & \ding{55} & GDA/GMM   & Dirichlet \\
    F-EDL~\citep{yoon2025fedl}      & \checkmark & \checkmark & $(p,\tau)$ + head SN & \ding{55} & Flexible Dirichlet \\
    \textbf{FI-EDL+SN (ours)}       & \checkmark & \checkmark & \ding{55} & \ding{55} & Dirichlet \\
    \bottomrule
  \end{tabular}
\end{table}

\subsection{Calibration and out-of-distribution detection}\label{sec:rw-calood}

\paragraph{Calibration.}
Guo \textit{et al.}~\citep{guo2017calibration} document the modern
miscalibration pattern -- deeper/wider networks with batch normalization
are systematically overconfident -- and propose post-hoc temperature
scaling as a near-optimal fix. We adopt their ECE estimator as our primary
calibration metric. Mukhoti \textit{et al.}~\citep{mukhoti2020focal} show
that training with focal loss already yields well-calibrated networks
without sacrificing accuracy; we cite this as the strongest train-time
calibration baseline outside the evidential family. More recent
train-time approaches, such as the class-adaptive label-smoothing margin
of CALS~\citep{liu2023cals}, push the calibration vs.\ accuracy
trade-off further but remain orthogonal to the evidential-Dirichlet
formulation studied here.

\paragraph{Post-hoc OOD detection.}
Hendrycks and Gimpel~\citep{hendrycks2017baseline} establish the maximum
softmax probability baseline and the AUROC/AUPR/FPR$95$ protocol that we
follow. ODIN~\citep{liang2018odin} combines temperature scaling and
gradient-direction input perturbations to improve MSP. Lee
\textit{et al.}~\citep{lee2018mahalanobis} fit class-conditional Gaussians
in the feature space and use Mahalanobis distance as the OOD score, a
direct ancestor of DAEDL's density-augmented evidential head. KNN-OOD
\citep{sun2022knn} dispenses with parametric assumptions and uses the
$k$-th-nearest-neighbor distance in normalized feature space. SNGP
\citep{liu2020sngp} is closest in spirit to our recipe: it combines
spectral normalization on the hidden layers with a Gaussian-process output
head to achieve distance-aware uncertainty in a single deterministic pass.
More recent post-hoc detectors operate purely on the trained network's
intermediate representations: ASH~\citep{djurisic2023ash} prunes and
rescales activations at inference, and GEN~\citep{liu2023gen} replaces
maximum softmax probability with a generalised-entropy score; both are
complementary to the training-time changes we study and could be layered
on top of our recipe at no retraining cost.

\paragraph{Spectral normalization and Lipschitz nets.}
Miyato \textit{et al.}~\citep{miyato2018sngan} introduce the
power-iteration-based spectral-norm constraint that we use; van Amersfoort
\textit{et al.}'s DUQ~\citep{vanamersfoort2020duq} is a deterministic
uncertainty model built around an RBF head and a two-sided gradient
penalty for bi-Lipschitz sensitivity; Bartlett
\textit{et al.}~\citep{bartlett2017spectrally} provide the spectrally-
normalized margin generalization bound that motivates the constraint
beyond OOD distance-awareness. The recipe of \autoref{sec:recipe}
intentionally adopts only the simplest of these mechanisms (SN on the
backbone, no GP head, no two-sided penalty) so that the contribution of
SN to a Dirichlet head can be isolated.

% =========================================================================
\section{Background}\label{sec:background}

\subsection{Dirichlet evidential classifiers}
Let $x\!\in\!\mathcal{X}$ be an input and $y\!\in\!\{1,\dots,K\}$ its class
label. Evidential deep learning (EDL) reinterprets the network output as a
Dirichlet distribution over the categorical class probabilities $p$ rather
than a single point estimate~\citep{sensoy2018edl}. Concretely, the backbone
$f_\theta$ maps $x$ to non-negative evidence $e(x)\!=\!\mathrm{ReLU}(f_\theta(x))
\!\in\!\mathbb{R}_{\ge 0}^{K}$ and the Dirichlet concentration parameters
are obtained by adding a unit prior,
\begin{equation}
  \alpha_c(x) \;=\; e_c(x) + 1, \qquad c=1,\dots,K,
  \label{eq:alpha}
\end{equation}
so that $\alpha\!\in\!\mathbb{R}_{>0}^{K}$ and $\alpha_0(x)\!=\!\sum_{c=1}^{K}
\alpha_c(x)$ is the total Dirichlet strength. The mean class probabilities
are $\hat p_c = \alpha_c/\alpha_0$, and the Sensoy~\textit{et~al.} formulation
uses the expected sum-of-squared-errors (Brier-like) loss
\begin{equation}
  \mathcal{L}_{\text{fit}}(\alpha,y) \;=\; \sum_{c=1}^{K}\Bigl[(y_c-\hat p_c)^2
  + \tfrac{\hat p_c(1-\hat p_c)}{\alpha_0+1}\Bigr],
  \label{eq:fit}
\end{equation}
with $y$ a one-hot label.

\subsection{KL-to-uniform regularizer}
To suppress evidence on incorrect classes, standard EDL adds a KL divergence
between a target-masked concentration
$\tilde\alpha = \alpha\odot(1-y)+y$ and the uniform Dirichlet
$\mathrm{Dir}(\mathbf{1})$~\citep{sensoy2018edl}:
\begin{equation}
  \mathcal{L}_{\text{kl}}(\alpha,y) \;=\;
  \mathrm{KL}\!\bigl(\mathrm{Dir}(\tilde\alpha)\,\|\,\mathrm{Dir}(\mathbf{1})\bigr).
  \label{eq:kl}
\end{equation}
The target-class masking $\alpha\odot(1-y)+y$ protects the correct class from
being driven toward $1/K$. Variants such as I-EDL, R-EDL and Re-EDL inherit
this annealed KL term with different scheduling and prior choices.

\subsection{Fisher information of a Dirichlet}
\label{sec:fim}
For $p\sim\mathrm{Dir}(\alpha)$ with $\alpha_0=\sum_c\alpha_c$, the Fisher
information matrix is
\begin{equation}
  I(\alpha)_{cc'} \;=\; \delta_{cc'}\,\psi'(\alpha_c) - \psi'(\alpha_0),
  \label{eq:fim}
\end{equation}
where $\psi'(\cdot)$ is the trigamma function. The trace, used by F-EDL and
by the FI-EDL controller as an adaptive scalar summary of the predicted
distribution's local curvature, has the closed form
\begin{equation}
  \mathrm{tr}\bigl(I(\alpha)\bigr) \;=\;
  \sum_{c=1}^{K}\psi'(\alpha_c) \;-\; K\,\psi'(\alpha_0).
  \label{eq:fimtrace}
\end{equation}
Two observations matter for what follows. First, $\mathrm{tr}(I(\alpha))$ is
dominated by the $\sum_c\psi'(\alpha_c)$ term: when evidence is moderate
($\alpha_c\!\approx\!2\text{--}3$, $\alpha_0\!\approx\!K\!\cdot\!2$) each
$\psi'(\alpha_c)\!\approx\!1/\alpha_c$ contributes $\mathcal{O}(1/\alpha_c)$
while the subtracted $K\psi'(\alpha_0)\!\approx\!K/\alpha_0\!\approx\!1/2$
remains $\mathcal{O}(1)$, so the trace scales as
$\mathrm{tr}(I(\alpha))\!=\!\mathcal{O}(K)$ across the typical operating
range. Second, since the trigamma is monotone decreasing, the trace varies
slowly with $\alpha$ — \autoref{sec:diagnosis} will show that the per-sample
standard deviation of $\mathrm{tr}(I(\alpha))$ on CIFAR-10 is only $\sim 6\%$
of its mean, which has direct consequences for any FIM-based adaptive
controller.

% =========================================================================
\section{Method}\label{sec:method}

\subsection{Diagnosing FI-EDL: the gate is saturated off}\label{sec:diagnosis}
The FI-EDL training objective combines the Brier-like fit
$\mathcal{L}_{\text{fit}}$ from \eqref{eq:fit} with a gated KL regularizer,
\begin{equation}
  \mathcal{L}_{\text{FI-EDL}} \;=\; \mathcal{L}_{\text{fit}}(\alpha,y)
    \;+\; \lambda(\alpha)\,\mathcal{L}_{\text{kl}}(\alpha,y),
    \qquad
  \lambda(\alpha) \;=\; \beta\,\exp\!\bigl(-\gamma\,\mathrm{tr}(I(\alpha))\bigr),
  \label{eq:fiedl}
\end{equation}
with $\beta\!=\!\gamma\!=\!1$ in the original formulation. The intent is that
$\lambda(\alpha)$ adapts the KL strength to the local sharpness of the
predicted Dirichlet: confident (high-trace) regions should be regularized
less, uncertain regions more. We instrumented this gate during the matched
re-training of FI-EDL on MNIST and CIFAR-10 and recorded the actual values
attained on the training set; see \autoref{tab:gate}.

\begin{table}[t]
  \centering
  \caption{Empirical values of the FI-EDL gate
    $\lambda(\alpha)=\exp(-\mathrm{tr}(I(\alpha)))$ during training. The
    effective KL contribution $\lambda\!\cdot\!\mathcal{L}_{\text{kl}}$ is
    five to seven orders of magnitude smaller than the fit term.}
  \label{tab:gate}
  \small
  \begin{tabular}{lccccc}
    \toprule
    & $\mathrm{tr}(I(\alpha))$ & $\overline{\lambda}$ & $\lambda_{\max}$
    & $\mathcal{L}_{\text{kl}}$ (raw) & $\mathcal{L}_{\text{fit}}$ \\
    \midrule
    MNIST    & 14.76 & $4.2\!\times\!10^{-7}$ & $1.8\!\times\!10^{-6}$
             & 63.1 & 0.025 \\
    CIFAR-10 & 14.48 & $3.8\!\times\!10^{-4}$ & $9.3\!\times\!10^{-3}$
             & 66.2 & 0.037 \\
    \bottomrule
  \end{tabular}
\end{table}

Three facts emerge. First, the Fisher trace sits at
$\mathrm{tr}(I(\alpha))\!\approx\!14.5\!\approx\!K$ throughout training for
$K\!=\!10$, exactly the $\mathcal{O}(K)$ scale predicted by
\eqref{eq:fimtrace} in \autoref{sec:fim}. A naive point evaluation gives
$\exp(-14.5)\!\approx\!5\!\times\!10^{-7}$, but because the per-sample
trace has non-trivial variance (standard deviation $0.23$ on MNIST and
$0.94$ on CIFAR-10, around a mean of $\sim\!14.5$), Jensen's inequality
lifts the empirical $\overline{\lambda}$ to $4.2\!\times\!10^{-7}$ on
MNIST and the substantially larger $3.8\!\times\!10^{-4}$ on CIFAR-10
(\autoref{tab:gate}). Second, multiplying by this gate yields an
effective KL contribution
$\overline{\lambda}\,\mathcal{L}_{\text{kl}}\!\approx\!2.6\!\times\!10^{-5}$
on MNIST, three orders of magnitude below the fit term
($\mathcal{L}_{\text{fit}}\!=\!0.025$), but
$\overline{\lambda}\,\mathcal{L}_{\text{kl}}\!\approx\!2.5\!\times\!10^{-2}$
on CIFAR-10, which is roughly $70\%$ of
$\mathcal{L}_{\text{fit}}\!=\!0.037$. The gate is therefore saturated
near zero only on MNIST; on CIFAR-10 it acts as a near-constant
\emph{moderate} regulariser, not the inert one that the MNIST headline
suggests. Third, the trace's per-sample variance is tiny relative to its
mean ($6.5\%$ on CIFAR-10, $1.6\%$ on MNIST), so the gate cannot deliver
the per-sample adaptivity the controller was designed for on either
dataset --- every sample sees essentially the same KL weight.

\paragraph{Implication.} The FIM controller fails its design goal on
both datasets, but the failure mode differs. On MNIST the gate saturates
and FI-EDL trains essentially on the Brier-like fit term alone. On
CIFAR-10 the gate is non-trivial but constant -- it acts as a uniformly
applied moderate KL coefficient rather than the per-sample adaptive
regulariser the original work proposes. In neither case does the
controller deliver adaptivity, not because of an implementation bug but
because the constants $\beta,\gamma$ were not calibrated against the
$\mathcal{O}(K)$ magnitude of the Fisher trace and because the trace
itself is too sample-invariant to drive any adaptive gate. The strong
calibration results reported for FI-EDL therefore reflect the underlying
fit term plus, on CIFAR-10, an accidentally well-tuned scalar KL
strength -- not an adaptive Fisher-information controller. A different
explanation is required, which we develop next.

\subsection{Proposed recipe}\label{sec:recipe}
Given that the FIM gate is inert (\autoref{sec:diagnosis}) yet the underlying
FI-EDL training is empirically well-calibrated, we ask what \emph{actually}
drives the calibration and OOD behaviour in the FI-EDL family and whether a
simpler design can match or surpass it. We propose three components:

\begin{enumerate}
  \item[\textbf{(i)}] \textbf{KL-free Brier fit.} Train with the fit term
    \eqref{eq:fit} alone, dropping the KL-to-uniform regularizer
    \eqref{eq:kl} entirely. This makes explicit what FI-EDL was already doing
    under saturation (\autoref{sec:diagnosis}) and removes a regularizer whose
    full-$\alpha$ form actively degrades calibration when re-enabled
    (\autoref{sec:why}).
  \item[\textbf{(ii)}] \textbf{Backbone spectral normalization.} Apply
    spectral normalization~\citep{miyato2018sngan} to every convolutional and
    linear layer of the backbone $f_\theta$ (not to the evidential output
    head). This enforces a bi-Lipschitz feature map and prevents the
    feature/logit blow-up that softmax-trained classifiers exhibit on far-OOD
    inputs, matching the rationale used by SNGP and DUQ
    ~\citep{liu2020sngp,vanamersfoort2020duq}. We apply spectral
    normalization only in natural-image settings; small-backbone settings
    such as MNIST/ConvNet pay a calibration tax under SN (see
    \autoref{sec:abl-sn} for the regression).
  \item[\textbf{(iii)}] \textbf{Total-evidence ($\alpha_0$) OOD score.} For
    OOD detection, score an input by its total Dirichlet strength
    $s(x)=\alpha_0(x)=\sum_c \alpha_c(x)$, with higher $s$ meaning more
    in-distribution. This is monotone-equivalent to the vacuity score
    $K/\alpha_0$ used elsewhere~\citep{sensoy2018edl} but with
    confidence-like orientation, which makes it a drop-in replacement for the
    maximum-softmax-probability score~\citep{hendrycks2017baseline} on a
    Dirichlet head.
\end{enumerate}

\begin{algorithm}[t]
  \caption{Spectral-normalized, KL-free EDL (training and inference).}
  \label{alg:recipe}
  \small
  \KwIn{labeled set $\mathcal{D}=\{(x_i,y_i)\}$, backbone $f_\theta$,
    evidence head $g_\phi$.}
  \KwSty{Setting:} natural-image regime --- apply spectral
    normalization to every conv/linear layer of $f_\theta$; for small
    backbones (e.g.\ MNIST ConvNet), omit SN to avoid the calibration
    regression documented in \autoref{sec:abl-sn}.\\
  \textbf{Training pass.}\\
  \quad For each minibatch: compute $\alpha = g_\phi(f_\theta(x))+1$ and
       update $\theta,\phi$ by minimizing $\mathcal{L}_{\text{fit}}(\alpha,y)$
       in \eqref{eq:fit}. No KL term.\\[2pt]
  \textbf{Inference pass.}\\
  \quad Classification: $\hat y(x) = \arg\max_c \alpha_c(x)$.\\
  \quad OOD score: $s(x) = \alpha_0(x) = \sum_c \alpha_c(x)$
       (higher = more in-distribution).\\
  \quad Misclassification score: $\max_c \hat p_c(x) = \max_c \alpha_c(x)/\alpha_0(x)$.
\end{algorithm}

The recipe adds no learnable parameters over standard EDL and requires no
auxiliary density model. The only architectural change is a spectral
normalization wrapper around existing backbone layers, supported natively in
PyTorch.

\subsection{Why each component}\label{sec:why}

\paragraph{(i) Dropping KL-to-uniform.}
The KL term \eqref{eq:kl} in its standard target-masked form
$\tilde\alpha=\alpha\odot(1-y)+y$ is designed to suppress evidence on
incorrect classes while leaving the correct class untouched. The FI-EDL
formulation, however, uses an unmasked variant that applies
$\mathrm{KL}(\mathrm{Dir}(\alpha)\,\|\,\mathrm{Dir}(\mathbf{1}))$ to the full
concentration vector, pulling \emph{every} $\alpha_c$ — including the target
class — toward~$1$ and the predicted probabilities toward $1/K$. Under the
saturated gate (\autoref{sec:diagnosis}) this is harmless, but when the
controller is re-enabled by Fisher-trace normalization or KL warmup, it
becomes the dominant force on the loss. We measure the resulting collapse
quantitatively in \autoref{sec:abl-gate}: CIFAR-10 ECE rises from
$5.19\%$ to $71.62\%$ even though accuracy is essentially unchanged
($89.95\%$ vs.\ $90.32\%$), exactly the severe-underconfidence signature.
A target-masked KL avoids the collapse but, as \autoref{sec:abl-masked} shows,
yields no metric improvement over baseline either, because the Fisher trace's
near-zero per-sample variance leaves the controller with no adaptive signal.
We therefore drop the KL term outright.

\paragraph{(ii) Spectral normalization on the backbone.}
Spectral normalization rescales each weight matrix $W$ to
$W/\sigma(W)$, where $\sigma(W)$ is its largest singular
value~\citep{miyato2018sngan}. Applied to every conv/linear layer of the
backbone, this enforces a global Lipschitz constant of one (per layer) and,
in turn, prevents the feature norm from blowing up as the input moves away
from the training distribution. The same mechanism underlies the
distance-awareness of SNGP~\citep{liu2020sngp} and DUQ
~\citep{vanamersfoort2020duq}: a bi-Lipschitz feature map ensures that
far-OOD inputs land in a feature region where the trained classifier head
\emph{can} produce low-confidence outputs, instead of being driven to
high-confidence by uncontrolled feature magnitudes. Empirically this
improves all five CIFAR-10 metrics over FI-EDL without SN
(\autoref{sec:abl-sn}): the same component that gives SNGP its OOD
robustness gives a Dirichlet head a usable $\alpha_0$ signal.

\paragraph{(iii) $\alpha_0$ as the OOD score.}
The maximum predicted probability $\max_c\hat p_c$ is the de-facto OOD
baseline~\citep{hendrycks2017baseline} but, on a Dirichlet head, it discards
the very quantity the head was designed to expose: the total evidence
$\alpha_0$. Vacuity $K/\alpha_0$ is the canonical evidential uncertainty
score~\citep{sensoy2018edl}, and $\alpha_0$ is its monotone inverse, so the
two yield identical AUROC/AUPR/FPR95. The practical difference is purely
orientation — $\alpha_0$ reads "higher means more in-distribution" — and
this matches the convention of every other confidence score, which avoids
sign-flip bugs in downstream tooling. The empirical payoff is large: on
CIFAR-10$\to$SVHN, switching the FI-EDL baseline (no SN) from
$\max_c\hat p_c$ to $\alpha_0$ moves the AUROC from $89.89$ to $91.72$
with no retraining; on CIFAR-10$\to$CIFAR-100 the same score swap on the
FI-EDL baseline moves AUROC from $83.80$ to $84.67$. The recipe's
improvements are therefore partly \emph{score selection}, not just
training-time changes.

\paragraph{Why does the minimal recipe beat a more elaborate baseline?}
Both Flexible EDL \citep{yoon2025fedl} and our recipe apply backbone
spectral normalization and drop the unmasked KL-to-uniform regulariser,
yet on CIFAR-10/VGG-16 we observe a $-1.89$ percentage-point ECE gap
($p\!<\!0.001$) and a $+2.17$ percentage-point CIFAR-100 AUROC gap
($p\!<\!0.01$) in favour of the simpler design (\autoref{sec:main},
\autoref{sec:abl-sn}). Two mechanisms are consistent with this
observation. First, F-EDL routes the prediction through three
sub-networks (concentration $\alpha$, allocation $p$, and dispersion
$\tau$) sharing the SN-constrained backbone; the $p$ and $\tau$ heads
inject additional gradient signal into shared features, and the
Flexible-Dirichlet effective concentration
$\alpha_{\text{eff}}\!=\!\alpha\!+\!\tau\!\cdot\!p$ splits the total
evidence across a mixture component, so the same numerical $\alpha_0$
does not carry the same ``total evidence'' semantics as in a plain
Dirichlet --- consistent with our observation (\autoref{sec:score-ablation})
that F-EDL's $\alpha_0$ score under-performs its $\max_c\hat p_c$ score
on CIFAR-10$\to$SVHN. Second, F-EDL's Brier-on-mean regulariser
$\|\mathbf{y}-\mathbf{p}\|^2$ duplicates information already present in
the fit term \eqref{eq:fit}, doubling the gradient pressure on the same
target signal and constraining the backbone harder than the Lipschitz
bound alone; the minimal recipe avoids this interference. We propose
these as hypotheses rather than verified causal mechanisms; a controlled
ablation comparing F-EDL with and without each component is left to
future work. The empirical takeaway holds either way: on the studied
benchmarks, adding the Flexible-Dirichlet machinery on top of an
SN-regularised Dirichlet head does not improve --- and on three of four
headline metrics measurably degrades --- the resulting calibration and
OOD performance.

% =========================================================================
\section{Experiments}\label{sec:exp}

\subsection{Setup}\label{sec:setup}

\paragraph{Datasets.}
We evaluate on two in-distribution (ID) image classification benchmarks:
MNIST and CIFAR-10. For each ID set we use two near/far OOD pairs:
MNIST is paired with KMNIST (Kuzushiji-MNIST, near OOD, same image domain)
and FMNIST (Fashion-MNIST, semantic near-OOD); CIFAR-10 is paired with
SVHN (digit photographs, distribution shift in style) and CIFAR-100
(natural-image semantic near-OOD with disjoint class set). All four OOD sets
are evaluated against the trained ID classifier without any OOD exposure at
training time.

\paragraph{Backbones.}
For MNIST we use a small ConvNet matching the architecture of prior EDL work
~\citep{sensoy2018edl}; for CIFAR-10 we use VGG-16 as the main backbone in
\autoref{tab:main-cifar10}, \autoref{tab:main-ood} and
\autoref{tab:main-misc}, and also report a parallel ResNet-18 evaluation
in \autoref{tab:main-resnet} to confirm backbone generality of the
recipe.

\paragraph{Baselines.}
We compare against six published Dirichlet/evidential methods: standard EDL
~\citep{sensoy2018edl}, I-EDL, R-EDL, Re-EDL, F-EDL, and DAEDL
~\citep{yoon2024daedl}. All baselines are re-trained under our matched
protocol using their published hyperparameters where prescribed; we do not
use the published numbers directly.

\paragraph{Training protocol.}
We use a matched protocol across every method for fairness. Both datasets
train for at most 200 epochs with Adam, batch size 64, and no learning-rate
scheduler. The learning rate is $1\!\times\!10^{-3}$ for MNIST/ConvNet and
$5\!\times\!10^{-4}$ for CIFAR-10/VGG-16. We hold out a $20\%$ split of the
training set on MNIST and a $5\%$ split on CIFAR-10 as validation, and use
\texttt{val/accuracy} with mode \texttt{max} as the early-stopping monitor
(patience 20 epochs). We use \texttt{val/accuracy} rather than
\texttt{val/loss} because the KL warmup schedule of several EDL variants
makes the validation loss non-monotone over the first 10 epochs and
incompatible with patience-based stopping. All experiments use 5 random
seeds $\{0,1,2,3,4\}$ in 32-bit precision.

\paragraph{Statistical procedure.}
We report mean and standard deviation across 5 paired seeds for every
metric. For the headline CIFAR-10 comparisons of FI-EDL+SN against the three
strongest competitors (DAEDL, F-EDL, Re-EDL) we additionally run paired
$t$-tests and 10\,000-resample bootstrap $95\%$ confidence intervals on the
per-seed differences, with each method scored using its own best
confidence/OOD score (DAEDL and Re-EDL at $\alpha_0$ for OOD, F-EDL at
$\max_c\hat p_c$, all methods at $\max_c\hat p_c$ for misclassification
detection). Full per-test numbers are in \autoref{sec:main} and in the
supplementary statistics file the supplementary material.
We adopt the convention that a difference is reported as a "win" only when
$p\!<\!0.05$ \emph{and} the bootstrap $95\%$ CI excludes zero, and as a
statistical "tie" otherwise; this convention is used uniformly in
\autoref{sec:main} and throughout the discussion.

\paragraph{Code and data availability.}
All datasets used in this work are publicly available standard
benchmarks: MNIST and CIFAR-10 as the in-distribution training sets,
KMNIST and FashionMNIST as MNIST OOD, and SVHN and CIFAR-100 as CIFAR-10
OOD. The implementation, Hydra configurations for every preset, training
scripts (\texttt{scripts/run\_*.sh}), evaluation scripts, and aggregated
results CSVs will be released at
\url{https://github.com/<anonymous>/FI-EDL} upon acceptance; an
anonymised snapshot is available to reviewers on request. All
experiments were run on a single NVIDIA GPU; per-seed wall-clock times
are approximately 15--20 minutes for MNIST (ConvNet), 1--1.5 hours for
CIFAR-10/VGG-16, and 50 minutes for CIFAR-10/ResNet-18, with no
specialised hardware required. Framework versions are pinned in the
released \texttt{pyproject.toml} (PyTorch 2.x, PyTorch Lightning 2.x,
Hydra 1.3).

\subsection{Main results}\label{sec:main}

We report main results in three tables: calibration (ECE, NLL,
accuracy, AURC) in \autoref{tab:main-cifar10} and \autoref{tab:main-mnist};
OOD detection (AUROC at each method's best score) in
\autoref{tab:main-ood}; and misclassification detection (AUROC at
$\max_c\hat p_c$) in \autoref{tab:main-misc}. The "FI-EDL" row throughout
refers to the recommended recipe of \autoref{sec:recipe}: on CIFAR-10 the
backbone uses spectral normalisation; on MNIST it does not, per the
SN-off prescription discussed in \autoref{sec:disc-sn} and ablated in
\autoref{sec:abl-sn}. All numbers are 5-seed mean $\pm$ standard
deviation; bold marks the unique statistically-significant best (paired
$t$-test $p\!<\!0.05$ \emph{and} bootstrap 95\% CI excludes zero, see
the supplementary material).

\begin{table*}[t]
  \centering
  \caption{Calibration and accuracy on CIFAR-10 with VGG-16 (5 seeds).
    ECE and accuracy in percent; NLL is negative log-likelihood; AURC is
    area under the risk-coverage curve. Bold = unique
    statistically-significant best; underline = best-mean but
    paired-tied. \textbf{FI-EDL row} is the recommended recipe
    (backbone spectral normalisation = on). Statistical significance is
    against the best competitor on each metric; full paired tests in the
    supplementary material.}
  \label{tab:main-cifar10}
  \small
  \begin{tabular}{lcccc}
    \toprule
    Method & ECE (\%)$\downarrow$ & NLL$\downarrow$ & Accuracy (\%)$\uparrow$ & AURC$\downarrow$ \\
    \midrule
    EDL    & $14.36 \pm 6.46$ & $1.6035 \pm 0.1361$ & $49.89 \pm 9.83$ & $0.2589 \pm 0.1380$ \\
    I-EDL  & $39.44 \pm 1.11$ & $0.8357 \pm 0.0233$ & $90.58 \pm 0.60$ & $0.0173 \pm 0.0014$ \\
    R-EDL  & $5.27 \pm 0.48$  & $0.4451 \pm 0.0264$ & $89.66 \pm 0.46$ & $0.0293 \pm 0.0030$ \\
    Re-EDL ($\lambda{=}0.8$) & $5.11 \pm 1.11$ & $0.4256 \pm 0.0308$ & $90.38 \pm 0.26$ & $0.0168 \pm 0.0011$ \\
    DAEDL  & $4.17 \pm 0.46$  & $0.4366 \pm 0.0106$ & $89.99 \pm 0.21$ & $0.0146 \pm 0.0011$ \\
    F-EDL  & $5.82 \pm 0.55$  & $0.4536 \pm 0.0438$ & $90.53 \pm 0.34$ & $0.0153 \pm 0.0006$ \\
    \textbf{FI-EDL (ours)} & \underline{$3.93 \pm 0.45$} & $0.3980 \pm 0.0040$ & $90.48 \pm 0.39$ & $0.0149 \pm 0.0008$ \\
    \bottomrule
  \end{tabular}
\end{table*}

\begin{table}[t]
  \centering
  \caption{Calibration and accuracy on MNIST with the small ConvNet
    (5 seeds). FI-EDL row uses the SN-off configuration of
    \autoref{sec:abl-sn}.}
  \label{tab:main-mnist}
  \small
  \begin{tabular}{lcccc}
    \toprule
    Method & ECE (\%)$\downarrow$ & NLL$\downarrow$ & Acc.\ (\%)$\uparrow$ & AURC$\downarrow$ \\
    \midrule
    EDL    & $6.73 \pm 1.78$  & $0.1281 \pm 0.0306$ & $98.34 \pm 0.31$ & $0.0003 \pm 0.0001$ \\
    I-EDL  & $37.68 \pm 0.25$ & $0.4978 \pm 0.0062$ & $99.49 \pm 0.06$ & $0.0001 \pm 0.0000$ \\
    R-EDL  & $5.97 \pm 1.12$  & $0.0820 \pm 0.0143$ & $99.38 \pm 0.07$ & $0.0001 \pm 0.0000$ \\
    Re-EDL & $9.33 \pm 1.11$  & $0.1179 \pm 0.0132$ & $99.44 \pm 0.08$ & $0.0001 \pm 0.0000$ \\
    DAEDL  & $50.97 \pm 1.63$ & $0.8933 \pm 0.0399$ & $95.83 \pm 0.30$ & $0.0027 \pm 0.0003$ \\
    F-EDL  & $\mathbf{0.48 \pm 0.30}$ & $0.0253 \pm 0.0041$ & $99.34 \pm 0.10$ & $0.0001 \pm 0.0001$ \\
    \textbf{FI-EDL (ours, SN-off)} & $1.70 \pm 0.73$ & $0.0846 \pm 0.1062$ & $97.31 \pm 4.44$ & $0.0013 \pm 0.0026$ \\
    \bottomrule
  \end{tabular}
\end{table}

\begin{table*}[t]
  \centering
  \caption{Out-of-distribution detection AUROC ($\uparrow$) on the four
    OOD pairs. Each method is scored at its own best OOD score
    (\autoref{sec:score-ablation}), shown in the ``Score'' column.
    Identical AUROC values for $\alpha_0$ and vacuity are reported as
    $\alpha_0$. CIFAR-10 OOD numbers are at VGG-16. Bold = unique
    statistically-significant best in the column; underline = best-mean
    but paired-tied.}
  \label{tab:main-ood}
  \small
  \begin{tabular}{llcccc}
    \toprule
    Method & Score & MNIST$\to$KMNIST & MNIST$\to$FMNIST & CIFAR-10$\to$SVHN & CIFAR-10$\to$C100 \\
    \midrule
    EDL    & $\max_c\hat p_c$ & $95.38 \pm 1.33$ & $91.04 \pm 2.45$ & $75.01 \pm 2.55$ & $62.79 \pm 3.03$ \\
    I-EDL  & $\max_c\hat p_c$ & $96.45 \pm 0.92$ & $92.36 \pm 3.48$ & $88.00 \pm 3.96$ & $82.41 \pm 1.23$ \\
    R-EDL  & $\max_c\hat p_c$ & $95.63 \pm 0.55$ & $95.43 \pm 0.95$ & $87.12 \pm 1.56$ & $79.92 \pm 0.62$ \\
    Re-EDL & $\alpha_0$       & $96.40 \pm 0.69$ & $96.08 \pm 1.15$ & $91.47 \pm 2.02$ & $85.82 \pm 0.30$ \\
    DAEDL  & $\alpha_0$       & $\mathbf{99.98 \pm 0.01}$ & $\mathbf{99.99 \pm 0.00}$ & $88.63 \pm 1.86$ & $84.68 \pm 0.45$ \\
    F-EDL  & $\max_c\hat p_c$ & $88.44 \pm 3.07$ & $97.22 \pm 0.81$ & $90.89 \pm 2.54$ & $84.15 \pm 0.71$ \\
    \textbf{FI-EDL (ours)} & $\alpha_0$ & $94.08 \pm 4.26$ & $90.44 \pm 6.83$ & \underline{$92.99 \pm 1.77$} & $\mathbf{86.32 \pm 0.20}$ \\
    \bottomrule
  \end{tabular}
\end{table*}

\begin{table}[t]
  \centering
  \caption{Misclassification detection on CIFAR-10/VGG-16 (5 seeds).
    AUROC ($\uparrow$) of $\max_c\hat p_c$ as a correct-vs-wrong score
    on the ID test set. Bold = unique statistically-significant best.}
  \label{tab:main-misc}
  \small
  \begin{tabular}{lcc}
    \toprule
    Method & AUROC$\uparrow$ & Accuracy (\%)$\uparrow$ \\
    \midrule
    EDL    & $81.01 \pm 8.96$ & $49.89 \pm 9.83$ \\
    I-EDL  & $89.39 \pm 0.75$ & $90.58 \pm 0.60$ \\
    R-EDL  & $76.03 \pm 2.72$ & $89.66 \pm 0.46$ \\
    Re-EDL & $89.40 \pm 0.60$ & $90.38 \pm 0.26$ \\
    DAEDL  & $\mathbf{91.99 \pm 0.53}$ & $89.99 \pm 0.21$ \\
    F-EDL  & $69.42 \pm 7.40$ & $90.53 \pm 0.34$ \\
    \textbf{FI-EDL (ours)} & $90.70 \pm 0.20$ & $90.48 \pm 0.39$ \\
    \bottomrule
  \end{tabular}
\end{table}

\paragraph{Headline reading of the tables.} On CIFAR-10/VGG-16, the
recipe attains the unique statistically-significant best on
CIFAR-10$\to$CIFAR-100 OOD AUROC ($86.32$, $p\!<\!0.01$ over each of
DAEDL, F-EDL, Re-EDL), best-mean ECE ($3.93$, paired-tied with DAEDL
$4.17$, $p\!=\!0.55$), best-mean SVHN AUROC ($92.99$, paired-tied with
Re-EDL $91.47$ at $p\!=\!0.12$ and F-EDL $90.89$ at $p\!=\!0.09$), and
second on misclassification AUROC ($90.70$ versus DAEDL $91.99$,
$p\!<\!0.01$). Accuracy at $90.48$ is statistically indistinguishable
from the top three baselines. On MNIST the recommended SN-off
configuration is second on ECE ($1.70$, behind F-EDL $0.48$), and
sixth-to-seventh on OOD AUROC where DAEDL's GMM density head is
empirically uncatchable by training-time-only signals
(\autoref{sec:disc-limits} (ii)).

\paragraph{Backbone generality: ResNet-18.}
To confirm the recipe is not VGG-16-specific, we repeat the matched
re-training of the four strongest contenders --- FI-EDL+SN, DAEDL,
F-EDL, Re-EDL --- with a ResNet-18 backbone on the same protocol (5
seeds, lr $5\!\times\!10^{-4}$, 200-epoch budget with val/accuracy
early stopping). Results appear in \autoref{tab:main-resnet}; full
paired-test statistics are in the supplementary material.

\begin{table*}[t]
  \centering
  \caption{CIFAR-10 with ResNet-18 (5 seeds). Backbone generality
    confirmation for the recipe. Bold marks the unique
    statistically-significant best per metric; underline marks
    best-mean with at least one paired-tie. Paired tests are
    FI-EDL+SN-resnet against each baseline-resnet.}
  \label{tab:main-resnet}
  \small
  \begin{tabular}{lccccc}
    \toprule
    Method & Acc (\%)$\uparrow$ & ECE (\%)$\downarrow$ & SVHN AUROC$\uparrow$
      & C100 AUROC$\uparrow$ & Misclass AUROC$\uparrow$ \\
    \midrule
    DAEDL  & $90.28 \pm 0.50$ & $3.73 \pm 0.26$ & $87.94 \pm 1.44$
      & $84.63 \pm 0.49$ & $\mathbf{92.93 \pm 0.58}$ \\
    F-EDL  & $91.71 \pm 0.19$ & $4.77 \pm 0.26$ & $85.85 \pm 1.25$
      & $84.63 \pm 0.75$ & $91.03 \pm 0.50$ \\
    Re-EDL & $91.83 \pm 0.32$ & $12.82 \pm 2.14$ & $92.51 \pm 1.20$
      & $\mathbf{87.94 \pm 0.39}$ & $91.31 \pm 0.50$ \\
    \textbf{FI-EDL (ours)} & $91.77 \pm 0.26$ & $\mathbf{3.15 \pm 0.57}$
      & \underline{$93.02 \pm 1.56$} & $87.08 \pm 0.40$ & $90.75 \pm 0.25$ \\
    \bottomrule
  \end{tabular}
\end{table*}

The ResNet-18 transfer reorders the headlines slightly. The recipe now
attains the \emph{uniquely significant} best ECE on the new backbone
($3.15$ vs DAEDL $3.73$, $p\!=\!0.038$; vs F-EDL $4.77$, $p\!=\!0.007$;
vs Re-EDL $12.82$, $p\!<\!0.001$) --- a stronger calibration claim than
on VGG-16, where DAEDL was a paired-tied competitor. SVHN OOD AUROC
remains best-mean ($93.02$), uniquely significant over DAEDL and F-EDL
but paired-tied with Re-EDL ($p\!=\!0.67$). On CIFAR-10$\to$CIFAR-100,
however, Re-EDL on ResNet-18 attains $87.94$ versus our $87.08$, a
statistically significant $-0.86$ percentage-point gap ($p\!<\!0.01$);
this is the only main-table metric on which the recipe is not at the
top of the ResNet-18 ranking. The Misclass-AUROC loss to DAEDL of
\autoref{tab:main-misc} is preserved, at a larger $-2.18$ percentage
points ($p\!<\!0.001$). Accuracy is statistically tied with F-EDL and
Re-EDL at the top ($\sim 91.8\%$). Re-EDL's exceptional ECE regression
on ResNet-18 ($12.82$, versus $5.11$ on VGG-16) is worth noting --- the
KL-free Re-EDL is calibrated only under a backbone with strong feature
regularisation, suggesting that an explicit Lipschitz constraint (our
recipe) is more robust to backbone choice than evidence-weighted prior
relaxation alone.

\paragraph{Cross-backbone summary.} Across VGG-16
(\autoref{tab:main-cifar10}, \autoref{tab:main-ood},
\autoref{tab:main-misc}) and ResNet-18 (\autoref{tab:main-resnet}), the
recipe is the only method that is best-mean or best-significant on
\emph{both} ECE and one CIFAR-10 OOD direction simultaneously. Its
weak spots are consistent: misclassification AUROC trails DAEDL on
both backbones (the missing density signal), and on ResNet-18
specifically CIFAR-100 OOD goes to Re-EDL by a narrow but significant
margin. These are documented honestly in \autoref{sec:disc-limits}.

\paragraph{Additional OOD sets: DTD and GTSRB.}
To check that the recipe's OOD wins do not depend on the specific
choice of CIFAR-10$\to$SVHN/CIFAR-100, we additionally evaluate on two
extended Hendrycks-style far-OOD benchmarks: Describable Textures
(DTD)~\citep{cimpoi2014describing} with $1{,}880$ textured natural
images, and the German Traffic Sign Recognition Benchmark
(GTSRB)~\citep{stallkamp2012gtsrb} with $12{,}630$ traffic-sign images.
GTSRB is one of the six OOD pairs used by Re-EDL in their CIFAR-10
setup~\citep{chen2025reedl}. Both are forward-only on the existing
VGG-16 checkpoints (5 seeds, each method at its best OOD score).
Results in \autoref{tab:main-extended-ood}.

On DTD the recipe attains $\mathrm{AUROC}\!=\!91.09\!\pm\!0.28$,
\emph{uniquely significantly} above each of DAEDL ($89.42$,
$p\!=\!0.017$), F-EDL ($87.95$, $p\!=\!0.003$), and Re-EDL ($90.00$,
$p\!=\!0.009$). On GTSRB it attains $\mathrm{AUROC}\!=\!88.82\!\pm\!1.30$,
significantly beating DAEDL ($86.61$, $p\!=\!0.001$) and F-EDL
($86.07$, $p\!=\!0.017$), and statistically tied with Re-EDL
($87.40$, $p\!=\!0.11$). Together with CIFAR-10$\to$CIFAR-100
(\autoref{tab:main-ood}), the recipe has three CIFAR-10 OOD sets where
it is uniquely significantly best, and two more (SVHN, GTSRB) where it
is best-mean and paired-tied with Re-EDL.

\begin{table}[t]
  \centering
  \caption{Extended CIFAR-10 OOD AUROC ($\uparrow$) on Describable
    Textures (DTD) and GTSRB (German Traffic Sign Recognition
    Benchmark), 5 seeds, VGG-16, each method at its best OOD score.
    Bold = unique statistically-significant best ($p\!<\!0.05$ and
    bootstrap CI excludes zero); underline = best-mean with at least
    one paired-tie.}
  \label{tab:main-extended-ood}
  \small
  \begin{tabular}{lcccc}
    \toprule
    Method & DTD AUROC ($\uparrow$) & DTD score & GTSRB AUROC ($\uparrow$) & GTSRB score \\
    \midrule
    F-EDL  & $87.95 \pm 0.80$ & $\max_c\hat p_c$ & $86.07 \pm 2.12$ & $\max_c\hat p_c$ \\
    DAEDL  & $89.42 \pm 0.74$ & $\alpha_0$ & $86.61 \pm 1.49$ & $\alpha_0$ \\
    Re-EDL & $90.00 \pm 0.47$ & $\alpha_0$ & $87.40 \pm 1.22$ & $\alpha_0$ \\
    \textbf{FI-EDL (ours)} & $\mathbf{91.09 \pm 0.28}$ & $\alpha_0$ & \underline{$88.82 \pm 1.30$} & $\alpha_0$ \\
    \bottomrule
  \end{tabular}
\end{table}

\paragraph{Harder benchmark: CIFAR-100 (F-EDL setup).}
To test whether the recipe transfers to a $10\times$ larger label space,
we replicate the F-EDL CIFAR-100 setup~\citep{yoon2025fedl}: ResNet-18
backbone, matched protocol (5 seeds, lr $5\!\times\!10^{-4}$, batch 64,
val/accuracy early stopping with patience 20), and OOD evaluation on
SVHN, CIFAR-10, TinyImageNet, and DTD. Results in
\autoref{tab:main-cifar100}. Three findings stand out. First, with the
larger $K\!=\!100$ the recipe attains the best mean accuracy
($67.78\!\pm\!0.55\%$, paired-tied with Re-EDL $68.21$) and best mean
ECE among methods that actually train successfully ($8.43\!\pm\!1.69\%$,
paired-tied with F-EDL $10.23$, $p\!=\!0.34$; significantly better than
Re-EDL's $29.17\%$, $p\!<\!0.001$). Second, on every OOD direction the
recipe is best-mean or paired-tied with Re-EDL: TIN
$75.95\!\pm\!0.56$ (paired-tied with Re-EDL $76.38$, $p\!=\!0.20$),
CIFAR-10 $73.57\!\pm\!1.20$ (Re-EDL $73.85$, $p\!=\!0.72$), SVHN
$75.02\!\pm\!3.51$ (Re-EDL $75.55$, $p\!=\!0.85$), DTD
$71.96\!\pm\!1.94$ (Re-EDL $70.65$, $p\!=\!0.27$). All four are
statistically tied at the top.

Third, DAEDL underperforms severely on this benchmark. During training
the network reaches validation accuracy $\sim\!20\%$ (well above
random $1/K\!=\!1\%$ but well below the $66\!\pm\!67\%$ of the other
three methods); the post-hoc class-conditional GMM is then fit
successfully on $47{,}500$ samples $\times\,512$-dimensional features
(per our training-time log), but the test-time density-weighted logit
$\alpha\!=\!\exp(g_{\phi}(f_{\theta}(x))\!\cdot\!s(x))$ further
compresses predicted probabilities toward uniform, dropping test
accuracy to $1.76\!\pm\!0.55\%$. F-EDL's published paper independently
reports DAEDL on the same benchmark at $66.01\!\pm\!2.6\%$ accuracy
but with ECE $86.00\%$~\citep{yoon2025fedl}, indicating that DAEDL on
CIFAR-100 either fails to train (our matched protocol) or trains to
reasonable accuracy but with severely under-confident predictions
(F-EDL's report). In either regime, DAEDL's density-augmented EDL
recipe is empirically brittle at $K\!=\!100$ in a way that the
density-free recipes (FI-EDL, Re-EDL, F-EDL) are not.

\begin{table*}[t]
  \centering
  \caption{CIFAR-100 with ResNet-18 (matching F-EDL Appendix E.2,
    5 seeds). Each method at its best score for the OOD column; ECE in
    percent; bold = unique statistically-significant best; underline =
    best-mean with at least one paired-tie. DAEDL's collapse to
    near-random accuracy is described in the surrounding text and
    matches the analogous degradation reported by \citet{yoon2025fedl}
    on the same benchmark.}
  \label{tab:main-cifar100}
  \small
  \begin{tabular}{lcccccc}
    \toprule
    Method & Acc (\%)$\uparrow$ & ECE (\%)$\downarrow$ & SVHN AUROC$\uparrow$
      & C10 AUROC$\uparrow$ & TIN AUROC$\uparrow$ & DTD AUROC$\uparrow$ \\
    \midrule
    DAEDL$^{\dagger}$  & $1.76 \pm 0.55$ & $0.76 \pm 0.55$ &
      $48.68 \pm 3.74$ & $49.93 \pm 0.53$ & $50.02 \pm 0.42$ & $46.88 \pm 1.13$ \\
    F-EDL   & $66.61 \pm 0.33$ & $10.23 \pm 3.33$ &
      $64.30 \pm 3.11$ & $70.55 \pm 0.36$ & $73.81 \pm 0.54$ & $65.09 \pm 1.67$ \\
    Re-EDL  & $\underline{68.21 \pm 0.50}$ & $29.17 \pm 3.60$ &
      \underline{$75.55 \pm 2.84$} & \underline{$73.85 \pm 0.50$} &
      \underline{$76.38 \pm 0.14$} & $70.65 \pm 0.72$ \\
    \textbf{FI-EDL (ours)} & $67.78 \pm 0.55$ & $\underline{8.43 \pm 1.69}$ &
      $75.02 \pm 3.51$ & $73.57 \pm 1.20$ & $75.95 \pm 0.56$ &
      $\underline{71.96 \pm 1.94}$ \\
    \bottomrule
  \end{tabular}
  \\ {\footnotesize $^{\dagger}$DAEDL on CIFAR-100/ResNet-18 reaches
   validation accuracy $\sim\!20\%$ during training (above random $1\%$
   but well below the other three methods); the post-hoc density
   weighting at test time then compresses predictions toward uniform,
   dropping test accuracy to $1.76\!\pm\!0.55\%$. F-EDL's
   paper~\citep{yoon2025fedl} reports DAEDL at $66\%$ accuracy with
   $86\%$ ECE on the same benchmark — an analogous, calibration-side
   collapse. See main text for the full mechanism.}
\end{table*}

\paragraph{Robustness under corruption: CIFAR-10-C
(DAEDL/F-EDL setup).}
We additionally evaluate distribution-shift robustness on
CIFAR-10-C~\citep{hendrycks2019benchmarking}, the
$19\!\times\!5\!=\!95$ corrupted variants of the CIFAR-10 test set used
as the standard robustness benchmark in both DAEDL (Table 4) and
F-EDL (Table 2). Forward-only on the existing VGG-16 checkpoints
(5 seeds), per-severity numbers averaged over the 19 corruption types
are in \autoref{tab:main-cifar10c}. Three findings, all paired-tested
across the 5 matched seeds:

(i) Classification accuracy under corruption.  The recipe is best at
every severity. Averaged across severities 1--5, it improves accuracy
by $+3.18$ percentage points over DAEDL ($p\!<\!0.001$), $+1.38$ over
F-EDL ($p\!=\!0.017$), and $+1.65$ over Re-EDL ($p\!=\!0.030$);
the gap widens with severity, reaching $+5.12$ over DAEDL at the
worst level $\mathcal{C}\!=\!5$ ($p\!<\!0.001$).

(ii) Calibration under corruption.  Averaged across severities, the
recipe attains ECE $16.81\%$ versus F-EDL $19.72$ ($p\!=\!0.027$,
WIN), Re-EDL $19.52$ ($p\!=\!0.14$, paired-tie), and DAEDL $14.43$
($p\!=\!0.042$, loss).  DAEDL's post-hoc class-conditional density
head, while empirically unstable at $K\!=\!100$
(\autoref{tab:main-cifar100}), gives it a small but significant
calibration advantage under continuous distribution shift on
CIFAR-10. This is one of two CIFAR-10 metrics on which the recipe
trails DAEDL --- the other being clean-test misclassification AUROC
(\autoref{tab:main-misc}) --- and motivates the future-work direction
of layering a density head on top of our recipe
(\autoref{sec:disc-future}).

(iii) Distribution-shift detection AUPR.  All four methods are
statistically tied at every severity: the maximum predicted
probability signal differentiates clean from corrupted inputs with
essentially identical AUPR ($\sim\!56$--$73\%$ across severities), so
the recipe's improvements over DAEDL/F-EDL/Re-EDL on cleaner OOD
benchmarks do not translate into a measurable shift-detection
advantage on CIFAR-10-C.

\begin{table*}[t]
  \centering
  \caption{CIFAR-10-C per-severity results, averaged over 19
    corruption types and 5 seeds.  Each cell is mean $\pm$ std across
    seeds (the inner 19-corruption average is deterministic per
    seed). Bold = unique statistically-significant best at that
    severity (paired $t$-test vs.\ second-best, $p\!<\!0.05$);
    underline = best-mean with at least one paired-tie.}
  \label{tab:main-cifar10c}
  \small
  \begin{tabular}{lccccc}
    \toprule
    Method & $\mathcal{C}\!=\!1$ & $\mathcal{C}\!=\!2$ & $\mathcal{C}\!=\!3$ & $\mathcal{C}\!=\!4$ & $\mathcal{C}\!=\!5$ \\
    \midrule
    \multicolumn{6}{l}{\emph{Classification accuracy under corruption (\%, higher is better)}} \\
    DAEDL  & $82.73 \pm 0.30$ & $76.61 \pm 0.65$ & $70.59 \pm 0.72$ & $63.15 \pm 0.73$ & $51.88 \pm 0.88$ \\
    F-EDL  & $83.82 \pm 0.47$ & $78.02 \pm 0.79$ & $72.36 \pm 1.08$ & $65.29 \pm 1.39$ & $54.47 \pm 1.64$ \\
    Re-EDL & $83.30 \pm 0.60$ & $77.45 \pm 0.95$ & $72.04 \pm 1.17$ & $65.25 \pm 1.39$ & $54.58 \pm 1.99$ \\
    \textbf{FI-EDL (ours)} & $\mathbf{84.07 \pm 0.34}$ & $\mathbf{78.67 \pm 0.45}$ & $\mathbf{73.71 \pm 0.62}$ & $\mathbf{67.39 \pm 0.76}$ & \underline{$57.00 \pm 1.22$} \\
    \midrule
    \multicolumn{6}{l}{\emph{ECE under corruption (\%, lower is better)}} \\
    DAEDL  & $\mathbf{7.56 \pm 0.27}$ & $\mathbf{10.48 \pm 0.16}$ & $\mathbf{13.46 \pm 0.14}$ & $\mathbf{17.30 \pm 0.38}$ & $\mathbf{23.34 \pm 0.83}$ \\
    F-EDL  & $10.39 \pm 0.89$ & $14.47 \pm 1.25$ & $18.56 \pm 1.80$ & $23.74 \pm 2.30$ & $31.44 \pm 3.02$ \\
    Re-EDL & $10.05 \pm 1.69$ & $14.31 \pm 2.30$ & $18.35 \pm 2.85$ & $23.48 \pm 3.49$ & $31.43 \pm 4.88$ \\
    \textbf{FI-EDL (ours)} & $8.32 \pm 0.79$ & $12.17 \pm 1.13$ & $15.72 \pm 1.46$ & $20.21 \pm 1.78$ & $27.61 \pm 3.05$ \\
    \midrule
    \multicolumn{6}{l}{\emph{Distribution-shift detection AUPR (\%, higher is better)}} \\
    DAEDL  & $56.67 \pm 0.41$ & $61.49 \pm 0.62$ & $65.12 \pm 0.70$ & $68.80 \pm 0.85$ & $73.40 \pm 0.93$ \\
    F-EDL  & $56.67 \pm 0.23$ & $61.54 \pm 0.39$ & $65.19 \pm 0.38$ & $68.84 \pm 0.38$ & $73.44 \pm 0.45$ \\
    Re-EDL & $56.45 \pm 0.46$ & $61.00 \pm 0.59$ & $64.33 \pm 0.51$ & $67.84 \pm 0.37$ & $72.37 \pm 0.14$ \\
    FI-EDL (ours) & $56.56 \pm 0.48$ & $61.26 \pm 0.83$ & $64.82 \pm 1.05$ & $68.62 \pm 1.11$ & $73.46 \pm 1.24$ \\
    \bottomrule
  \end{tabular}
\end{table*}

\subsection{Score selection ablation}\label{sec:score-ablation}

Headline OOD comparisons in the literature typically report a single
confidence score (most often $\max_c\hat p_c$) across all methods, which
risks under-reporting evidential methods whose total evidence $\alpha_0$
is a stronger separator. We re-evaluate every baseline at four scores
$\{\max_c\hat p_c,\ \max_c\alpha_c,\ \alpha_0,\ \text{vacuity}\}$ on
CIFAR-10$\to$\{SVHN, CIFAR-100\}, with no retraining; the full matrix is
in the supplementary material and \autoref{tab:score-audit} summarises
the per-method best score.

\begin{table}[t]
  \centering
  \caption{Per-method best OOD score on CIFAR-10/VGG-16, 5-seed AUROC.
    \emph{vacuity} $=\!K/\alpha_0$ is a monotone-decreasing transform of
    $\alpha_0$, so the two yield identical ranking and identical AUROC;
    we group them and report $\alpha_0$ (higher = more in-distribution).
    Numbers are extracted from the same 5-seed runs that appear in
    \autoref{tab:main-cifar10}.}
  \label{tab:score-audit}
  \small
  \begin{tabular}{lccc}
    \toprule
    Method & SVHN best & CIFAR-100 best & $\alpha_0\!-\!\max_c\hat p_c$ swap (SVHN / C100) \\
    \midrule
    EDL    & $\max_c\hat p_c$ & $\max_c\hat p_c$ & $-6.17$ / $-1.79$ \\
    I-EDL  & $\max_c\hat p_c$ & $\max_c\hat p_c$ & $-0.50$ / $-0.47$ \\
    R-EDL  & $\max_c\hat p_c$ & $\max_c\hat p_c$ & $-25.44$ / $-9.44$ \\
    F-EDL  & $\max_c\hat p_c$ & $\max_c\hat p_c$ & $-54.70$ / $-26.87$ \\
    DAEDL  & either           & either           & $+0.02$ / $+0.01$ \\
    Re-EDL & $\alpha_0$       & $\alpha_0$       & $+1.57$ / $+1.22$ \\
    FI-EDL & $\alpha_0$       & $\alpha_0$       & $+1.83$ / $+0.87$ \\
    \bottomrule
  \end{tabular}
\end{table}

\paragraph{Three findings.}
First, the right OOD score depends on the loss family. Methods that
keep the target-class evidence numerically meaningful --- Re-EDL drops
the KL term so $\alpha_0$ tracks total evidence directly, and FI-EDL
trains without any KL crushing the target $\alpha$ --- consistently
prefer $\alpha_0$ over $\max_c\hat p_c$, with swap gains of $+1.2$ to
$+1.8$ percentage points. Methods that aggressively regularise toward
the uniform Dirichlet (EDL, I-EDL, R-EDL) or that split total evidence
across a mixture component (F-EDL's
$\alpha_{\text{eff}}\!=\!\alpha\!+\!\tau\!\cdot\!p$) instead degrade
catastrophically at $\alpha_0$ --- F-EDL's $\alpha_0$ AUROC drops to
$36.19$ on SVHN, $54.7$ percentage points below its $\max_c\hat p_c$
score. The same effect at smaller magnitude appears for R-EDL ($-25.4$
percentage points). Second, DAEDL is the unique baseline for which both
scores are within $0.02$ percentage points of each other: its post-hoc
GMM density head scales the predicted distribution so that both
$\max_c\hat p_c$ and $\alpha_0$ track the same underlying density
signal. Third, on misclassification detection (correct-vs-wrong on the
ID test set) a separate audit (not shown for brevity) confirms that
$\max_c\hat p_c$ remains the best score across every method including
ours; the $\alpha_0$ advantage is OOD-specific. We therefore report all
headline OOD numbers in \autoref{sec:main} at each method's best OOD
score and all misclassification numbers at $\max_c\hat p_c$. A single
fixed score across all methods is a measurement convention that
systematically misrepresents at least three baselines on the
\autoref{tab:main-cifar10} OOD comparison; this measurement caveat,
rather than a new model, is contribution 5 of \autoref{sec:intro}.

% =========================================================================
\section{Ablations}\label{sec:ablations}

\subsection{FIM controller: gate-on causes calibration
collapse}\label{sec:abl-gate}

\autoref{sec:diagnosis} establishes that under the published
$\beta\!=\!\gamma\!=\!1$ the FIM gate is saturated near zero, so the KL
regularizer is effectively absent during FI-EDL training. A natural
question is what happens once the gate is restored to a non-trivial value.
We test three configurations on CIFAR-10 with two seeds each, using
recent normalization tricks intended to bring $\mathrm{tr}(I(\alpha))$ to
$\mathcal{O}(1)$ before exponentiation: (i) a fixed division by the class
count $K$ (\texttt{div\_k}) so that the gate stabilizes at
$\lambda\!\approx\!0.23$; (ii) a per-batch $z$-score
(\texttt{batch\_z}) with $\gamma\!=\!0.5$; and (iii) the same per-batch
$z$-score but driven by the total evidence $\alpha_0$ rather than by the
Fisher trace. A standard 10-epoch KL warmup is applied in every case.
\autoref{tab:abl-gate-on} reports the result.

\begin{table}[t]
  \centering
  \caption{Stage-1 sweep on CIFAR-10/VGG-16. Re-enabling the FI-EDL gate
    with any of three signal-normalization schemes collapses calibration
    (ECE rises from $5.19\%$ to $71.62$--$78.67\%$) at essentially
    unchanged accuracy, the unmistakable signature of severe
    underconfidence. Baseline is FI-EDL with the gate left off
    ($\lambda\!\approx\!5\!\times\!10^{-7}$); 5 seeds for the baseline,
    2 seeds for the other rows.}
  \label{tab:abl-gate-on}
  \small
  \begin{tabular}{lcccc}
    \toprule
    Configuration & Acc (\%) $\uparrow$ & ECE (\%) $\downarrow$
      & SVHN AUROC $\uparrow$ & C100 AUROC $\uparrow$ \\
    \midrule
    Baseline (gate off)
      & $90.32\!\pm\!0.31$ & $\mathbf{5.19\!\pm\!0.54}$
      & $\mathbf{91.72\!\pm\!2.21}$ & $\mathbf{84.67\!\pm\!1.59}$ \\
    Fisher + div\_k, $\gamma\!=\!1$ (gate on)$^{\dagger}$
      & $89.95$ & $71.62$ & $92.31$ & $79.41$ \\
    Fisher + batch\_z, $\gamma\!=\!0.5\,^{\dagger}$
      & $89.75$ & $78.67$ & $88.27$ & $75.65$ \\
    $\alpha_0$ + batch\_z, $\gamma\!=\!0.5\,^{\dagger}$
      & $88.69$ & $77.29$ & $69.72$ & $68.42$ \\
    \bottomrule
  \end{tabular}
  \\ {\footnotesize $^{\dagger}$ $n\!=\!2$ seeds (diagnostic sweep); the
   baseline row uses $n\!=\!5$.}
\end{table}

The mechanism is immediate. FI-EDL's KL term uses the full Dirichlet
concentration $\mathrm{KL}(\mathrm{Dir}(\alpha)\,\|\,\mathrm{Dir}(\mathbf{1}))$,
\emph{including the target class}, rather than the target-masked form
$\tilde\alpha\!=\!\alpha\odot(1-y)+y$ used by standard EDL. With the gate
saturated at $5\!\times\!10^{-7}$ this is harmless; with the gate restored,
the KL term actively pulls every $\alpha_c$ -- including the predicted
class -- toward unity, and the predicted maximum probability collapses
toward $1/K\!=\!0.1$. Accuracy is preserved because the argmax remains
stable, but the calibration metric quantifies precisely this gap between
the (suppressed) confidence and the (unchanged) accuracy. Any meaningful
KL regularization in the original FI-EDL form therefore destroys
calibration. Conclusion: the gate is not a controller that can be
``unlocked'' to do useful work; the form of the regularizer itself is
incompatible with target-class confidence.

\subsection{Masked-KL controller does not save the
recipe}\label{sec:abl-masked}

The natural next attempt is to swap the full-$\alpha$ KL in FI-EDL for the
standard target-masked form $\tilde\alpha\!=\!\alpha\odot(1-y)+y$, which
protects the predicted class from being driven to $1/K$. If the
controller's adaptivity is meaningful, the masked variant should now
deliver a net gain over the gate-off baseline. We test the two surviving
configurations from \autoref{tab:abl-gate-on} -- the only ones whose
accuracy did not regress -- with the masked KL, two seeds each. Results
are in \autoref{tab:abl-masked}.

\begin{table}[t]
  \centering
  \caption{Stage-2 sweep on CIFAR-10/VGG-16. Replacing the full-$\alpha$
    KL with the target-masked variant $\tilde\alpha\!=\!\alpha\odot(1-y)+y$
    prevents the calibration collapse but yields no metric gain over the
    gate-off baseline. The $\alpha_0$-driven variant additionally fails
    to train.}
  \label{tab:abl-masked}
  \small
  \begin{tabular}{lcccc}
    \toprule
    Configuration & Acc (\%) $\uparrow$ & ECE (\%) $\downarrow$
      & SVHN AUROC $\uparrow$ & C100 AUROC $\uparrow$ \\
    \midrule
    Baseline (gate off)
      & $\mathbf{90.32\!\pm\!0.31}$ & $\mathbf{5.19\!\pm\!0.54}$
      & $\mathbf{91.72\!\pm\!2.21}$ & $\mathbf{84.67\!\pm\!1.59}$ \\
    Fisher + div\_k, $\gamma\!=\!1$ + masked$^{\dagger}$
      & $89.26\!\pm\!0.49$ & $5.46\!\pm\!0.82$
      & $90.86\!\pm\!0.40$ & $84.00\!\pm\!0.48$ \\
    $\alpha_0$ + batch\_z, $\gamma\!=\!0.5$ + masked$^{\dagger}$
      & $38.43\!\pm\!0.69$ & $9.44\!\pm\!5.44$
      & $63.70\!\pm\!5.63$ & $56.50\!\pm\!1.58$ \\
    \bottomrule
  \end{tabular}
  \\ {\footnotesize $^{\dagger}$ $n\!=\!2$ seeds (diagnostic sweep); the
   baseline row uses $n\!=\!5$.}
\end{table}

Masking recovers calibration ($5.46$ vs.\ $5.19\%$ baseline ECE) but
adds nothing: accuracy drops by $1.06$ pp, SVHN AUROC by $0.86$, CIFAR-100
AUROC by $0.67$. The root cause is the one already exposed by the
diagnosis in \autoref{sec:diagnosis}: the per-sample standard deviation of
$\mathrm{tr}(I(\alpha))$ is only $6.5\%$ of its mean, so even after
normalization the controller has essentially no per-sample variation to
adapt to. The masked-KL+\texttt{div\_k} configuration is therefore
empirically indistinguishable from a fixed-strength KL term, which the
EDL literature already explored. The $\alpha_0$-driven variant collapses
training entirely (38\% accuracy) because the controller signal is dominated
by the batch-level normalization rather than by sample-level evidence.

The combined conclusion of \autoref{sec:abl-gate} and \autoref{sec:abl-masked}
is that the FIM controller mechanism is empirically inert under the
original FI-EDL framing, that any straightforward rescue either preserves
the calibration tax of full-$\alpha$ KL or loses the adaptivity premise,
and that the cleanest path forward is to remove the controller entirely.
This is what the recipe of \autoref{sec:recipe} does.

\subsection{Backbone spectral normalization: CIFAR-10 gain, MNIST
regression}\label{sec:abl-sn}

The final ablation isolates the effect of backbone spectral normalization
in the proposed recipe. We re-train FI-EDL with the gate-off, KL-free fit
of \autoref{sec:recipe} on both MNIST (ConvNet) and CIFAR-10 (VGG-16),
toggling \texttt{backbone\_spectral\_norm} on or off and holding everything
else fixed. Five seeds per cell. Results are in
\autoref{tab:abl-sn}.

\begin{table}[t]
  \centering
  \caption{Backbone spectral normalization on/off, KL-free fit, 5 seeds.
    SN improves every CIFAR-10 metric but harms MNIST calibration:
    ECE jumps from $1.70\%$ to $13.60\%$, and the $\alpha_0$ OOD score
    becomes unusable on FMNIST. Reported scores: AUROCs use $\alpha_0$ on
    CIFAR-10 and the score listed on MNIST.}
  \label{tab:abl-sn}
  \small
  \begin{tabular}{llcccc}
    \toprule
    Dataset & Config & Acc (\%) $\uparrow$ & ECE (\%) $\downarrow$
      & OOD$_1$ AUROC $\uparrow$ & OOD$_2$ AUROC $\uparrow$ \\
    \midrule
    \multirow{3}{*}{MNIST}
      & baseline (SN off, maxp) & $97.31\!\pm\!4.44$
        & $\mathbf{1.70\!\pm\!0.73}$
        & KMNIST $\mathbf{94.08}$ & FMNIST $\mathbf{90.44}$ \\
      & +SN (maxp)              & $\mathbf{98.82\!\pm\!0.08}$
        & $13.60\!\pm\!1.27$
        & KMNIST $92.06$ & FMNIST $\mathbf{93.38}$ \\
      & +SN ($\alpha_0$)        & --- & --- & KMNIST $81.30$ &
        FMNIST $47.04$ \\
    \midrule
    \multirow{2}{*}{CIFAR-10}
      & baseline (SN off, $\alpha_0$) & $90.32\!\pm\!0.31$
        & $5.19\!\pm\!0.54$ & SVHN $91.72\!\pm\!2.21$
        & C100 $84.67\!\pm\!1.59$ \\
      & \textbf{+SN ($\alpha_0$)}     & $\mathbf{90.48\!\pm\!0.39}$
        & $\mathbf{3.93\!\pm\!0.45}$
        & SVHN $\mathbf{92.99\!\pm\!1.77}$
        & C100 $\mathbf{86.32\!\pm\!0.20}$ \\
    \bottomrule
  \end{tabular}
\end{table}

\paragraph{CIFAR-10: every metric improves.} SN reduces ECE from
$5.19$ to $3.93\%$, lifts CIFAR-10$\to$SVHN AUROC from $91.72$ to $92.99$,
and CIFAR-10$\to$CIFAR-100 AUROC from $84.67$ to $86.32$, at unchanged
accuracy. The result is consistent with the distance-aware Lipschitz
argument of SNGP and DUQ~\citep{liu2020sngp,vanamersfoort2020duq}:
constraining the spectral norm of every backbone layer prevents the
feature norm from blowing up on far-OOD inputs, so the Dirichlet head's
total evidence $\alpha_0$ remains a usable confidence signal across the
input space rather than being dominated by uncontrolled magnitudes.

\paragraph{MNIST: SN regresses calibration.} On the same recipe, SN
\emph{worsens} ECE from $1.70$ to $13.60\%$ and inverts the score
priority: the $\alpha_0$ score becomes unusable on FMNIST (AUROC
$47.04$, below chance) while $\max_c\hat p_c$ remains marginally usable.
Accuracy and seed-variance do improve ($97.31\!\pm\!4.44$ to
$98.82\!\pm\!0.08\%$), but the calibration cost is severe.
Mechanistically, the small MNIST ConvNet is already close to the
Lipschitz budget that SN enforces, so the additional constraint
over-tightens the logit magnitude: the network becomes systematically
underconfident, and total evidence $\alpha_0$ no longer scales with
input familiarity because it is pinned near a small value across the
whole input space. The same constraint that liberates a deep VGG-16 from
overconfident far-OOD features over-regularizes the small MNIST backbone.

\paragraph{Practical prescription.} The recipe is therefore SN-on for
natural-image backbones (where it sets a new CIFAR-10 OOD SOTA, see
\autoref{sec:main}) and SN-off for small-backbone settings such as MNIST.
This split is consistent with the SNGP/DUQ literature, where the
distance-awareness argument assumes a backbone deep enough that the
Lipschitz budget is not already at the bottleneck. A learned SN
coefficient that automatically adapts the constraint to the backbone is
a natural next step (\autoref{sec:discussion}).

% =========================================================================
\section{Discussion}\label{sec:discussion}

\subsection{Why does a deliberately minimal recipe match or beat more
elaborate baselines?}\label{sec:disc-why}
The skeptical reading of our results is: ``if you only \emph{remove}
machinery and still win, what is your method actually adding?'' Three
threads of evidence give a coherent answer. First, additional machinery
adds additional failure modes. F-EDL's allocation and dispersion
sub-heads $(p,\tau)$ ($\sim\!270$ thousand extra trainable parameters in
the CIFAR-10/VGG-16 configuration) introduce extra gradient pathways
into the shared backbone; DAEDL's post-hoc Gaussian density estimate
depends on class-conditional features being well separated; and the
original FI-EDL Fisher-information controller relies on an adaptive
signal that, as we showed in \autoref{sec:diagnosis}, has only
$6.5\%$ relative per-sample variation. Each additional component is
empirically inert or actively detrimental in at least one of our
benchmarks (\autoref{sec:abl-gate}, \autoref{sec:abl-masked},
\autoref{sec:abl-sn}). Second, score semantics matter. The audit in
\autoref{sec:score-ablation} shows that F-EDL, while built on a
Flexible-Dirichlet effective concentration
$\alpha_{\text{eff}}\!=\!\alpha\!+\!\tau\!\cdot\!p$, has an $\alpha_0$
AUROC of $36.19$ on SVHN --- $54.7$ percentage points below its
$\max_c\hat p_c$ score, indicating that the total-evidence quantity is
no longer monotone with confidence under the Flexible-Dirichlet design.
The plain Dirichlet head retains clean $\alpha_0$ semantics. Third, the
active ingredient of distance-aware Dirichlet uncertainty is the
bi-Lipschitz backbone, in line with the SNGP and DUQ analysis
\citep{liu2020sngp,vanamersfoort2020duq}, and that ingredient is shared
across our recipe, DAEDL, and F-EDL. Once the bi-Lipschitz backbone is
in place, adding more machinery on top of it does not appear to give an
additional Dirichlet-head signal --- and in the F-EDL case appears to
degrade it.

\subsection{When does spectral normalisation help?}\label{sec:disc-sn}
\autoref{sec:abl-sn} documents two regimes for the SN ablation. On
CIFAR-10 with VGG-16, applying SN to every convolutional and linear
layer of the backbone improves every headline metric simultaneously
(ECE, both OOD AUROCs, misclassification AUROC, accuracy). On MNIST
with the small ConvNet, the same SN protocol regresses calibration
from ECE $1.70$ to $13.60\%$ and inverts the orientation of the
$\alpha_0$ OOD score on FMNIST (AUROC $47.04$, below chance). The
mechanistic difference, suggested by the SNGP analysis, is the ratio of
the Lipschitz constraint to the data manifold's intrinsic complexity:
on a high-capacity backbone solving a natural-image task, the unit-norm
spectral budget per layer is loose enough that the backbone can still
learn discriminative features, and the bi-Lipschitz constraint suppresses
runaway feature magnitudes on OOD inputs. On a small ConvNet solving
MNIST, the same per-layer constraint is tight relative to the backbone's
capacity, so logit magnitudes are uniformly suppressed; the resulting
maximum predicted probabilities cluster around $1/K\!=\!0.1$, and
calibration breaks. The practical implication is that SN is an
architectural choice that should be ablated per dataset/backbone, not a
default. Our recipe therefore recommends SN-on for natural-image-scale
backbones (VGG-16, ResNet-18 in our experiments; presumably ResNet-50,
ViTs, EfficientNets in larger settings) and SN-off for small backbones
or small datasets. A learned per-layer spectral coefficient is a natural
follow-up that may close the MNIST regression without losing the
CIFAR-10 gain.

\subsection{The $\alpha_0$ measurement convention.}\label{sec:disc-score}
A side-effect of the score audit (\autoref{sec:score-ablation}) is a
methodological observation: the published evidential-uncertainty
literature is inconsistent about which score to use for OOD AUROC, and
the inconsistency materially affects leaderboards. For instance,
reporting every method at a single fixed score (whether $\max_c\hat p_c$
or $\alpha_0$) would systematically misrepresent at least three of the
baselines we evaluate. We therefore advocate the simple convention of
reporting OOD numbers for each method at its own best-AUROC score, with
the full per-score matrix in the supplementary material. This is a
measurement contribution rather than a model contribution; it does not
require a new training method, and it explains roughly half of the gap
between our headline numbers and prior published numbers for FI-EDL on
the same benchmarks.

\subsection{Limitations.}\label{sec:disc-limits}
We acknowledge five specific limitations of the present work.
(i)~\textbf{MNIST regression under SN.} On the small-backbone MNIST
benchmark, applying our recipe with SN regresses ECE from $1.70$ to
$13.60\%$ (\autoref{sec:abl-sn}); we therefore recommend SN-off in that
regime, but a deployment without per-task ablation would silently
under-report calibration. (ii)~\textbf{Misclassification loss.} On
CIFAR-10 misclassification AUROC, our recipe trails DAEDL by a
statistically significant $1.29$ percentage points
($p\!<\!0.01$, \autoref{sec:main}); we attribute this to the
post-hoc density head, which gives DAEDL an explicit feature-distance
signal that our minimal recipe does not produce.
(iii)~\textbf{Backbone family.} CIFAR-10 results cover VGG-16 and
ResNet-18; MNIST covers only ConvNet. The recipe has not been studied on
ViT, ResNet-50, or EfficientNet variants, and we cannot rule out
backbone-specific failure modes at those scales. (iv)~\textbf{Scale.}
CIFAR-10 is our largest natural-image ID benchmark; ImageNet-100 and
ImageNet-1k results are absent for compute-budget reasons. The
bi-Lipschitz backbone argument used by SNGP and DUQ extends in principle
to ImageNet, and our recipe would simply layer Brier-only fit and
$\alpha_0$ scoring on top, but this is unverified. (v)~\textbf{Modality
and task.} We study image classification only. Evidential regression
\citep{amini2020evidential} and structured-output evidential methods are
not covered, and the simplification principle may or may not transfer.

\subsection{Future work.}\label{sec:disc-future}
Three directions follow naturally. First, a per-layer learned spectral
coefficient that interpolates between the unconstrained and unit-norm
extremes may close the MNIST calibration regression while preserving the
CIFAR-10 gains. Second, a hybrid that combines our minimal training-time
recipe with a DAEDL-style post-hoc density head may recover the
misclassification AUROC currently lost to DAEDL while keeping the
CIFAR-10 calibration and OOD wins; this is conceptually a clean
addition because the two contributions are orthogonal. Third, scaling to
ImageNet-100 / ImageNet-1k with ResNet-50 or ViT backbones would test
whether the recipe's CIFAR-10 advantage is a small-data artifact or a
genuine architectural finding. We also note that the CIFAR-10-C
corruption benchmark \citep{hendrycks2017baseline} would be a natural
robustness extension, especially given that the bi-Lipschitz backbone
should in principle stabilise calibration under corruption; preliminary
infrastructure for this evaluation is already prepared in our codebase.

% =========================================================================
\section{Conclusion}\label{sec:conclusion}

We revisited the FI-EDL family of evidential uncertainty methods and
asked what is empirically responsible for its calibration and OOD
strength. The originally proposed Fisher-information controller turns
out to be empirically inert: the trace $\mathrm{tr}\,I(\alpha)\!=\!
\mathcal{O}(K)$ throughout training, the gate
$\lambda\!=\!\beta\,\exp(-\gamma\,\mathrm{tr}\,I(\alpha))$ saturates to a
near-constant scalar (negligible on MNIST, $\sim\!70\%$ of the fit term
on CIFAR-10), and the per-sample variance of the trace is too small
($6.5\%$ relative standard deviation on CIFAR-10, $1.6\%$ on MNIST) for
any adaptive controller to do real work. We then identified a
deliberately minimal three-component recipe --- KL-free Brier-like fit,
backbone spectral normalisation for natural-image scale, and total
evidence $\alpha_0$ as the OOD score --- that matches or exceeds the
strongest current EDL baselines on the studied CIFAR-10/VGG-16 protocol:
a uniquely statistically significant first place on CIFAR-10$\to$CIFAR-100
out-of-distribution detection (AUROC $86.32\!\pm\!0.20$,
$p\!<\!0.01$ against each of DAEDL, F-EDL, and Re-EDL), best-mean
calibration (ECE $3.93\!\pm\!0.45\%$, paired-tied with DAEDL,
$p\!=\!0.55$), and best-mean SVHN AUROC ($92.99\!\pm\!1.77$, paired-tied
with Re-EDL and F-EDL). The same recipe regresses MNIST calibration
under spectral normalisation, and we therefore recommend SN as a
task-conditional rather than default component. Taken together the
results support an Occam's-razor reading of recent evidential deep
learning: the active ingredients of strong calibration and OOD
performance are surprisingly few, and adding additional heads, density
estimators, or adaptive regularisers does not necessarily help. Per-layer
spectral tuning, hybrid density-aware variants, and scaling to ImageNet
remain open.

% =========================================================================
\bibliographystyle{elsarticle-num}
\bibliography{refs}

\end{document}
